\chapter{Definizione di Software}
\label{chap:software}
\section{Definizione}
Un \textbf{software} è un programma che speficica le istruzioni che un calcolatore deve eseguire al fine di raggiungere uno scopo. Nel contesto dell'Ingegneria del Software, si includono nella definizione anche tutti i documenti che descrivono lo stesso nelle varie fasi della produzione del sistema. \\
L'\textbf{Ingegneria del Software} è la disciplina che agisce globalmente su tutte le componenti software (inclusi i documenti) al fine di produrre software definibile "di qualità". 
\section{(Alcune) Classificazioni dei Software}
\subsection{Software di sistema}
I \textbf{Software di Sistema} sono programmi che gestiscono l'hardware di un computer e forniscono servizi di base per far funzionare altri software. \\ 
\textbf{ES:} Sistemi operativi, driver di dispositivo
\subsection{Software applicativo}
I \textbf{Software Applicativi} sono programmi usati dall'utente finale per svolgere compiti specifici. \\
\textbf{ES:} Microsoft Word, Google Chrome
\subsection{Software scientifico e tecnico}
I \textbf{Software Scientifici e Tecnici} sono progettati per applicazioni scientifiche, ingegneristiche o matematiche complesse. \\
\textbf{ES:} Matlab, AutoCAD
\subsection{Software Embedded}
I \textbf{Software Embedded} sono software incorporati in dispositivi hardare specifici, con funzionalità spesso limitate e progettati per uno scopo specifico. \\
\textbf{ES:} sistema operativo di una TV, firmware di una lavatrice
\subsection{Software di Produttività}
I \textbf{Software di Produttività} sono progettati per aiutare gli utenti a svolgere compiti legati al lavoro o alla gestione del tempo e delle risorse. \\
\textbf{ES:} Excel, Notion
\subsection{Applicazioni Web}
Software accessibili via browser, eseguiti da un server in remoto e che non richiedono installazione locale. \\
\textbf{ES:} Gmail, Facebook
\subsection{Software per l'Intelligenza Artificiale}
Software progettati per simulare capacità umane come apprendimento, ragionamento o comprensione del linguaggio. \\
\textbf{ES:} ChatGPT, Alexa
\section{4P dello sviluppo software}
\subsection{People - Persone}
Le persone coinvolte nel progetto o che ne traerebbero vantaggio. Definite anche "{stakeholder}" o "{attori}", possono essere a grandi linee divise in:
\begin{itemize}
    \item \textbf{Business Management} - dirigenti, responsabili aziendali o chi ha una visione strategia del progetto. Stabiliscono gli obiettivi, decidono budget e tempistiche, approvano o fermano progetti in base ai risultati.
    \item \textbf{Project Manager} - hanno la responsabilità di organizzare e supervisionare il progetto giorno per giorno. Coordinano le attività del team, gestiscono tempi, costi, rischi e qualità, producono report e roadmap e mediano tra le esigenze tecniche e le richieste del cliente.
    \item \textbf{Development Team} - sviluppatori, designer, tester, analisti che lavorano direttamente alla realizzazione del software. Sviluppano il codice, testano le funzionalità, risolvono problemi tecnici e partecipano alla progettazione tecnica.
    \item \textbf{Customers} - clienti, persone o aziende che commissionano il software o che pagano per averlo sviluppato. Definiscono requisiti e aspettative, forniscono input funzionali e valutano il prodotto.
    \item \textbf{End Users} - utenti finali, persone che realmente usano il software nel loro quotidiano. Utilizzano il prodotto in casi reali e possono fornire feedback.
\end{itemize}
\subsection{Product - Prodotto}
Il Prodotto è l'insieme di funzionalità e servizi che il team consegna agli utenti o ai clienti, nonché l'obiettivo finale dello sviluppo. 
\subsection{Project - Progetto}
Nel Progetto vengono specificate le attività da svolgere per realizzare il prodotto e i limiti a cui le stesse potrebbero incorrere. Fanno parte del progetto concetti come scope, budgeting, deadline, milestone e rischi.
\subsection{Process - Processo}
Il processo definisce quali sono e come organizzare le attività necessarie alla realizzazione di un software. Il Capitolo \ref{chap:processi} si concentrerà sulla specifiche di differenti prospettive per l'analisi dei processi e su alcuni modelli specifici.
\section{Ciclo di vita del software}
Il \textbf{ciclo di vita} di un software è l'insieme di stati che lo stesso attraversa dal suo concepimento iniziale fino alla sua dimissione. \\
Le fasi principali del ciclo di vita sono:
\begin{enumerate}
    \item \textbf{comunicazione}, identificazione iniziale delle funzionalità necessarie;
    \item \textbf{pianificazione}, valutazione dei rischi e dei costi;
    \item \textbf{modellazione}, progettazione delle funzionalità attraverso opportuni modelli;
    \item \textbf{costruzione}, implementazione delle scelte e verifica delle stesse;
    \item \textbf{deployment}, rilascio nell'ambiente;
    \item \textbf{manuntenzione}, operazioni post-rilascio, possono essere di tre tipi:
    \begin{enumerate}
        \item \textbf{correttiva}, correzione di bug ed errori tecnici;
        \item \textbf{adattiva}, adattamento del sistema a piattaforme o tecnologie differenti;
        \item \textbf{perfettiva}, introduzione di nuove funzionalità;
    \end{enumerate}
\end{enumerate}